\newpage




\begin{thebibliography}{99}
\makeatletter
\renewcommand{\@biblabel}[1]{}
\makeatother
\bibitem{}
\noindent Ci si limita a
citare solo i libri (prevalentemente in lingua italiana)  la cui
impostazione metodologica \`{e} simile a quella qui adottata. Si
citano dapprima dei testi universitari di \emph{Scienza delle
Costruzioni} dove la Meccanica della trave rettilinea costituisce una parte
significativa dell'intero volume. Successivamente si menzionano
altri testi, di natura  monografica, su cui \`{e} possibile
approfondire alcuni argomenti specifici. Si riporta infine un breve
elenco di testi universitari che contengono gli argomenti \emph{di
base} necessari alla comprensione del volume.
\newline


\subsubsection{ \emph{Testi universitari di
Scienza delle Costruzioni}}
 \noindent La teoria della trave  tradizionalmente oggetto
dei corsi universitari di Scienza delle Costruzioni \`{e} trattato
nei seguenti testi.

\makeatletter
\renewcommand{\@biblabel}[1]{[#1]}
\setcounter{enumiv}{0} \makeatother



\bibitem{Benvenuto}
E. Benvenuto {\em La Scienza delle Costruzioni e il suo sviluppo
storico}, Manuali Sansoni Firenze, 1981.


\bibitem{Capurso}
M. Capurso {\em Lezioni di Scienza delle Costruzioni}, Pitagora Editrice Bologna, 1971.


%\bibitem{Casini_Vasta} P. Casini, M. Vasta  \emph{Scienza delle Costruzioni}, CittàStudi Edizioni, 2011.



\bibitem{Ceradini3}
G. Ceradini {\em La teoria della trave}, in Scienza
delle Costruzioni a cura di G. Ceradini, Vol. 3, ESA Roma, 1991.


\bibitem{Ceradini4}
G. Ceradini {\em Teoria lineare dei sistemi di travi}, in Scienza
delle Costruzioni a cura di G. Ceradini, Vol. 4, ESA Roma, 1991.



\bibitem{Corradi1}
L. Corradi Dell'Acqua {\em Meccanica delle strutture: 2 
Le teorie strutturali ed il metodo degli elementi finiti }, McGraw-Hill, 1992.

\bibitem{Podio1}
P. Podio Guidugli {\em Lezioni di Scienza delle Costruzioni}, Aracne editrice, 2009.



\vspace{0.5 cm} e, pi\`{u} specificamente, alla  plasticit\`{a}:


%\bibitem{CapursoGav}
%M. Capurso, C. Gavarini {\em Elasto-plasticit\`{a} delle
%strutture} in Plasticit\`{a}, vol. II, a cura di R. Baldacci, G.
%Ceradini e E. Giangreco, Milano, Italsider (Tamburini), 1971.


\bibitem{Corradi2}
L. Corradi Dell'Acqua {\em Meccanica delle strutture: 3 La
valutazione della capacit\`{a} portante}, McGraw-Hill, 1994.

\bibitem{Maier}
G. Maier {\em Complementi di scienza delle costruzioni}, clup,
1976.


\vspace{0.6 cm}
\subsubsection{\emph{Argomenti di base}}

 \noindent Le tematiche  di base di Geometria e Analisi
 necessarie alla comprensione del volume sono, ad esempio, sviluppate in:


\bibitem{Abeasis}
S. Abeasis {\em Elementi di algebra lineare e geometria}, Zanichelli,  1993.


\bibitem{Borisenko}
A.I. Borisenko, I.E. Tarapov {\em Vector and tensor analysis with
applications}, Dover Publications, Inc., New York 1979.



\bibitem{bramanti}
M. Bramanti, C.D. Pagani, S. Salsa \emph{Matematica, calcolo
infinitesimale e algebra lineare}, Zanichelli, 2000.


\bibitem{Dicuonzo}
V. Dicuonzo {\em Introduzione all'algebra lineare e alla geometria analitica}, Veschi,  1991.



\bibitem{Gasparini}
I. C. Gasparini {\em Strutture algebriche e operatori lineari}, Eredi Virgilio Veschi,  1975.



\bibitem{Ghizzetti1}
A. Ghizzetti, F. Rosati {\em Analisi matematica}, vol. I e II,
MASSON editoriale Veschi, 1992.


\bibitem{Pagani}
C.D. Pagani, S. Salsa \emph{Analisi matematica} vol. 1 e 2, Masson, 1998.

\bibitem{Vaccaro}
G. Vaccaro, A. Carfagna, L. Piccolella  {\em Lezioni di geometria e  algebra lineare}, Zanichelli,  1995.

\vspace{0.5 cm} \noindent Nozioni di Meccanica  generale sono
trattate in:


\bibitem{Benvenuti}
P. Benvenuti, G. Maschio \emph{Appunti delle lezioni di meccanica
razionale}, Kappa Editore, Roma, 2000.


\bibitem{Bordoni}
P. G. Bordoni \emph{Lezioni di meccanica razionale}, Zanichelli,
1992.

\bibitem{LCA} E.N.M. Cirillo  \emph{Appunti delle lezioni di meccanica
razionale per ingegneria}, Edizioni CompoMat, 2018.

\bibitem{DenHartog2}
J. P. Den Hartog \emph{Mechanics}, Dover Pubblications, Inc., New
York, 1961.

% \bibitem{Carpinteri2}
% A. Carpinteri \emph{Structural Mechanics Fundamentals}, CRC Press., 2014.

% \bibitem{Carpinteri1}
% A. Carpinteri \emph{Structural Mechanics - A unified approach}, E \& FN SPON, 1997.

%\bibitem{LCA}
%T. Levi-Civita, U. Amaldi \emph{Lezioni di meccanica razionale},
%(3 vol.), Zanichelli, 1991.



\vspace{0.5 cm} \noindent e, pi\`{u} specificamente, nozioni di
Meccanica orientate alla Scienza delle costruzioni sono contenute
in:

\bibitem{DB} D. Bernardini \emph{Statica - Un'introduzione alla meccanica delle strutture}, CittàStudi Edizioni, 2009. 

\bibitem{CeradiniSCR}
G. Ceradini {\em Cinematica e statica dei sistemi rigidi}, in Scienza
delle Costruzioni a cura di G. Ceradini, Vol. 1, ESA Roma, 1991.


\bibitem{Corradi1A}
C. Comi, L. Corradi Dell'Acqua {\em Introduzione alla meccanica
strutturale}, McGraw-Hill, 2003.

\bibitem{GBGN}
E. Guagenti, F. Buccino, E. Garavaglia, G. Novati \emph{Statica,
fondamenti di meccanica strutturale}, McGraw-Hill, seconda edizione,
2004.


\bibitem{LP}
A. Luongo, A. Paolone {\em Meccanica delle strutture: sistemi
rigidi ad elasticit\`{a} concentrata}, CEA, seconda edizione,
2004.

\bibitem{LPDN}
A. Luongo, A. Paolone, S. Di~Nino~: \textit{Rigid Structures with Point-Flexibility}, Springer, 2025.


\bibitem{Podio2}
P. Podio Guidugli {\em Lezioni di statica}, Aracne editrice, 2014.


\vspace{0.5 cm} \noindent Una trattazione della Meccanica
dei solidi e del problema di De Saint Venant coerente con l'impostazione qui seguita \`{e}
contenuta in:

\bibitem{LP1}
A. Luongo, A. Paolone {\em Scienza delle costruzioni, 1. Il
continuo di Cauchy}, CEA, 2004.

\bibitem{LP2}
A. Luongo, A. Paolone {\em Scienza delle costruzioni, 2. Il
problema di De Saint Venant}, CEA, 2005.

\bibitem{PVV}
A. Paolone, F. Vestroni, S. Vidoli {\em Scienza delle costruzioni,  l'analisi della tensione nelle travi - un software applicativo }, CEA, 2009.

Infine, può essere utile cominciare a consultare qualche manuale tecnico quale:

\bibitem{Cremonese}
{\em Manuale di Ingegneria civile - sezione seconda}, Edizioni scientifiche A. Cremonese - Roma, 1982.



\end{thebibliography}
